\section{A radial deficit and separate endpoint laws}\label{sec:radial}
After normalizing the lens radius to one, this argument uses the
norm-one head-test sums $p(S_i)$ as replacement coefficients; the
joint-budget proof used $\|x\|p(S_i)$. Estimating the coefficient
change and squared-test slack separately gives a radial-deficit bound.
The same calculation then controls each endpoint law.

\begin{proposition}[A radial-deficit estimate]\label{thm:tail}
For a finite ancestral $H$, $x=P_Hx$, and arbitrary $y\in X$ with
$\|x\pm y\|\le R$, one has
\begin{equation}\label{eq:atomic-tail}
 \|Q_Hy\|\le(\sqrt2+\sqrt6)\sqrt{R(R-\|x\|)}.
\end{equation}
\end{proposition}

\begin{proof}
If $R=0$, both vectors vanish. If $x=0$, contractivity of $Q_H$ gives
$\|Q_Hy\|\le R$, which implies the theorem. In the remaining case,
scale to $R=1$ and put $d=1-\|x\|$. By
Lemma~\ref{lem:finite-hull}, choose a norming test $p=P_Hp$ with
$\rho(p)=1$ and $x(p)=\|x\|$.

Fix a finite ancestral $K\supset H$. Represent each of
$P_K(x+y)$ and $P_K(x-y)$ as a finite convex combination of atoms,
as supplied by Lemma~\ref{lem:finite-hull}. On a finite probability
space choose a uniform sign $\sigma\in\{-1,1\}$ and conditionally
choose a represented atom $A=\sum_i a_i d_{S_i}$ from the corresponding
combination. Then
\[
 \mathbb EA=x,\qquad\mathbb E(\sigma A)=P_Ky.
\]
For each outcome set $b_i=p(S_i)$, $A^\circ=\sum_i b_i d_{S_i}$,
and $L=1-\sum_i b_i^2$. Replacing the original coefficients by these
segment sums has the joint error estimate
\[
 L+\sum_i(a_i-b_i)^2
 =1+\sum_i a_i^2-2A(p)\le2(1-A(p)).
\]
Taking expectations and applying \eqref{eq:atom-coefficients} gives
\begin{equation}\label{eq:mixture-error}
 \mathbb EL\le2d,\qquad
 \mathbb E\|A-A^\circ\|\le\sqrt{2d}.
\end{equation}

We have paid for coefficient replacement. It remains to control the
tail of the replaced signed mean. At a gate $g$, let $b(g)$ be the
crossing sum of the represented family, or zero. Segments entirely
outside $H$ have replaced coefficient zero. Thus
$A^\circ(e_g)=b(g)$ and its restriction to $J(U_g)$ has norm at most
$|b(g)|$. Define finite gate vectors
\[
 h(g)=\mathbb Eb(g),\qquad m(g)=\mathbb E|b(g)|.
\]
Because $x(e_g)=0$, Lemma~\ref{lem:cones} and
\eqref{eq:mixture-error} give $\|h\|_2\le\sqrt{2d}$.
Take two independent copies of the atom data. Lemma~\ref{lem:opposite-crossing}
and independence yield
\[
 \begin{split}
 \|m\|_2^2-\|h\|_2^2
 &=\mathbb E\sum_g\bigl(|b(g)b'(g)|-b(g)b'(g)\bigr)\\
 &\le\mathbb E(L+L')\le4d.
 \end{split}
\]
Consequently $\|m\|_2\le\sqrt{6d}$. By the cone identity
\eqref{eq:cone-decomposition},
\[
 \|Q_H\mathbb E(\sigma A^\circ)\|\le\|m\|_2\le\sqrt{6d}.
\]
Its distance to $Q_HP_Ky$ is at most $\sqrt{2d}$ by
\eqref{eq:mixture-error}. Let $K$ increase through an ancestral
exhaustion. Lemma~\ref{lem:forest-structure} gives $P_Ky\to y$, proving
\eqref{eq:atomic-tail} for $R=1$. Scaling gives the stated theorem.
\end{proof}

\subsection{A maximum bound from the radial estimate}

The radial estimate also gives a bound for the maximum-based modulus
defined in \eqref{eq:max-modulus}, with a weaker constant than the
joint-budget estimate. We retain this calculation to record the effect
of the alternative normalization.

\begin{corollary}\label{cor:max-modulus}
For the specified dual norm on $X$,
\[
 \overline\theta_X(t)\ge\frac{\sqrt{1+t^2/4}-1}{2}\qquad(t>0).
\]
\end{corollary}

\begin{proof}
For a finite-head unit vector $u=P_Hu$ and unit $y\in\ker P_H$, put
$R=\max\{\|u+ty\|,\|u-ty\|\}$. Convexity gives $R\ge1$.
Since $(\sqrt2+\sqrt6)^2<16$, Proposition~\ref{thm:tail} implies
$t^2\le16R(R-1)$. Solving this quadratic inequality gives the asserted
lower bound for $R-1$.

For an arbitrary unit $x$, choose an ancestral head $H$ with
$\|Q_Hx\|<\eta<1$. Then $u=P_Hx/\|P_Hx\|$ is well defined
and $\|x-u\|<2\eta$. The preceding bound for $R-1$, diminished
by $2\eta$, therefore holds for all unit $y\in\ker P_H$ with center
$x$. Take the infimum over these $y$, the supremum over finite-codimensional
subspaces, and then let $\eta\downarrow0$. The bound is independent
of $x$, so the final infimum proves the corollary.
\end{proof}

\subsection{Keeping the endpoint laws separate}

The proof of Proposition~\ref{thm:tail} estimates a signed mixture of the
two atom laws. It does not require a small replacement error for either
conditional law separately. The next proposition supplies that
representation-level information.

\begin{proposition}\label{prop:endpoint}
Let $0\le d<1$, let $H$ be finite and ancestral, and suppose
$x=P_Hx$, $\|x\|\ge1-d$, and $\|x\pm y\|\le1$. Fix a head
norming test $p$ with $\rho(p)=1$ and $x(p)=\|x\|$, and a finite
ancestral set $K\supset H$. For each $\tau\in\{-1,1\}$ choose
any finite represented atom law with mean $P_K(x+\tau y)$, and denote
its expectation by $\mathbb E_\tau$. For an atom $A=\sum_i a_i d_{S_i}$,
put $b_i=p(S_i)$, $A^\circ=\sum_i b_i d_{S_i}$, and
$L=1-\sum_i b_i^2$. Then
\begin{equation}\label{eq:endpoint-coefficients}
 \mathbb E_\tau\!\left[L+\sum_i(a_i-b_i)^2\right]\le4d,
 \qquad \mathbb E_\tau\|A-A^\circ\|\le2\sqrt d.
\end{equation}
If $b(g)$ is the replaced crossing sum at gate $g$, put
$m_\tau(g)=\mathbb E_\tau|b(g)|$ and
$m=(m_1+m_{-1})/2$. These nonnegative vectors satisfy
\begin{equation}\label{eq:endpoint-gates}
 m_\tau\le2m\text{ coordinatewise},\qquad\|m\|_2\le\sqrt{6d}.
\end{equation}
\end{proposition}

\begin{proof}
The test exists by Lemma~\ref{lem:finite-hull}, since
$\|x\|\ge1-d>0$. Both endpoint evaluations on $p$ are at most one,
so $|y(p)|\le1-\|x\|\le d$ and $(x\pm y)(p)\ge1-2d$.
The coefficient identity from the proof of Proposition~\ref{thm:tail} gives
\[
 \mathbb E_\tau\!\left[L+\sum_i(a_i-b_i)^2\right]
 \le2\bigl(1-(x+\tau y)(p)\bigr)\le4d.
\]
Equation~\eqref{eq:atom-coefficients} and Cauchy--Schwarz give the
second bound in \eqref{eq:endpoint-coefficients}.

Combine the laws with equal probabilities. Their mean is $x$, whose
pairing with $p$ is at least $1-d$. The coefficient and opposite-crossing
calculation in the proof of Proposition~\ref{thm:tail} therefore gives
$\|m\|_2\le\sqrt{6d}$. The coordinate inequality in
\eqref{eq:endpoint-gates} follows from $m=(m_1+m_{-1})/2$.

These separate-law bounds also give a tail estimate directly. The cone
identity implies
\[
 \|Q_H\mathbb E_\tau A^\circ\|\le\|m_\tau\|_2\le2\sqrt{6d},
 \qquad Q_H\mathbb E_\tau A=\tau Q_HP_Ky.
\]
Thus \eqref{eq:endpoint-coefficients} gives
$\|Q_HP_Ky\|\le2\sqrt d+2\sqrt{6d}$. Ancestral exhaustion yields
\[
 \|Q_Hy\|\le(2+2\sqrt6)\sqrt d.
\]
The larger constant records the cost of retaining the two conditional
laws separately; Proposition~\ref{thm:tail} uses their signed mixture.
\end{proof}
